sections/runtime.tex

% ============================

% Runtime Model and Execution Strategy

% ============================

\\section{Runtime Model and Execution Strategy}

The Ω⁺ runtime integrates discrete graph evolution, pattern-directed

rewrites, and continuous field dynamics into a unified execution model.

This section defines the runtime phases, scheduling policies, and

numerical integration strategy used to evaluate Ω⁺ programs.

\\subsection{Runtime Configuration}

Execution proceeds over configurations:

\\\[

&nbsp;\\Config = (G, F, R, P, \\tau, Q, M)

\\\]

as defined in Section~4. The configuration evolves through a sequence of

small-step transitions and epoch-level updates.

\\subsection{Epoch Structure}

Execution is organized into \\emph{epochs}, each of duration \$\\Delta t\$.

An epoch consists of five ordered phases:

\\begin{enumerate}

&nbsp;\\item \\textbf{Scheduled Dynamics.}

&nbsp;  All dynamics \$(n @ t_d ::= n')\$ with \$t_d \\le \\tau\$ are dequeued

&nbsp;  and applied. These may update node labels, introduce new nodes, or

&nbsp;  schedule additional events.

&nbsp;\\item \\textbf{Structural Rewrites.}

&nbsp;  All patterns \$p \\in P\$ are matched against the graph \$G\$.

&nbsp;  A scheduler selects a non-overlapping subset of matches, and each

&nbsp;  selected match fires, applying its associated rewrite.

&nbsp;\\item \\textbf{Continuous Field Integration.}

&nbsp;  Each continuous field (e.g., \$I\$, \$\\Mu\$) is updated by numerical

&nbsp;  integration of its governing PDE over \$\\Delta t\$.

&nbsp;\\item \\textbf{Discrete Field Stabilization.}

&nbsp;  Discrete field updates (e.g., \$\\Mu := \\Mu + \\eta \\nabla(\\Phi)\$)

&nbsp;  are applied based on the updated graph and field configuration.

&nbsp;\\item \\textbf{Region and Metadata Update.}

&nbsp;  Region memberships are recomputed, derived quantities are updated,

&nbsp;  and the logical time is advanced:

&nbsp;  \\\[

&nbsp;    \\tau \\leftarrow \\tau + \\Delta t.

&nbsp;  \\\]

\\end{enumerate}

The epoch transition is:

\\\[

&nbsp;\\Config_\\tau \\LStep{\\text{epoch}(\\Delta t)} \\Config_{\\tau + \\Delta t}.

\\\]

\\subsection{Scheduling Policies}

Pattern firing is controlled by a scheduler specified in \$M\$.

Supported policies include:

\\begin{itemize}

&nbsp;\\item \\textbf{Maximal Parallelism.}

&nbsp;  All non-overlapping matches fire simultaneously.

&nbsp;\\item \\textbf{Priority-Based Scheduling.}

&nbsp;  Patterns are assigned priorities; higher-priority patterns fire

&nbsp;  first, and lower-priority patterns may be deferred.

&nbsp;\\item \\textbf{Randomized Scheduling.}

&nbsp;  A random subset of matches is selected, enabling nondeterministic

&nbsp;  exploration of system behavior.

&nbsp;\\item \\textbf{Deterministic Sequential Scheduling.}

&nbsp;  Matches are ordered deterministically (e.g., by node identifier)

&nbsp;  and fired sequentially.

\\end{itemize}

Schedulers may enforce fairness guarantees, such as weak fairness:

if a pattern remains continuously enabled, it must eventually fire.

\\subsection{Numerical Integration}

Continuous fields evolve according to PDEs such as:

\\\[

&nbsp;\\frac{\\partial I}{\\partial t} = \\Mu(x)\\nabla I + \\rho(x,t).

\\\]

The runtime supports multiple integration methods:

\\begin{itemize}

&nbsp;\\item explicit Euler,

&nbsp;\\item Runge--Kutta (e.g., RK4),

&nbsp;\\item implicit solvers for stiff systems,

&nbsp;\\item spectral or finite-volume methods for structured domains.

\\end{itemize}

The choice of solver and discretization parameters (e.g., spatial step

\$h\$, time step \$\\Delta t\$) is stored in \$M\$.

Stability constraints such as the CFL condition must be satisfied:

\\\[

&nbsp;\\Delta t \\le \\frac{h}{\\max_x |\\Mu(x)|}.

\\\]

\\subsection{Region Recalculation}

Regions may be defined by spatial domains, graph subsets, or contextual

conditions. After each epoch, region membership is recomputed based on:

\\begin{itemize}

&nbsp;\\item node coordinates or embeddings,

&nbsp;\\item field thresholds,

&nbsp;\\item structural properties of the graph.

\\end{itemize}

\\subsection{Execution Parameters}

The metadata \$M\$ contains:

\\begin{itemize}

&nbsp;\\item scheduler choice,

&nbsp;\\item time step \$\\Delta t\$,

&nbsp;\\item field solver parameters,

&nbsp;\\item learning rates for discrete updates,

&nbsp;\\item numerical clamps and bounds,

&nbsp;\\item random seed for nondeterministic scheduling.

\\end{itemize}

These parameters allow Ω⁺ to be executed deterministically or

nondeterministically, depending on the desired behavior.

This runtime model provides a coherent framework for integrating

graph-based rewrites and continuous field evolution within a single

execution semantics.

